% Experiment Template
% 1) Copy this file and rename it, e.g. experiments/experiment_04_firewall_rules.tex
% 2) In main.tex, add:
%    \includeexperiment{XX}{Experiment Title}{experiments/your_new_file_name}
% 3) Keep this file content as sections only (no \experiment command here)

\section{Theory}
\subsection*{Hill Cipher}
The Hill cipher, proposed by Lester S. Hill, is a \textbf{polyalphabetic block cipher}.
Unlike stream ciphers (which process one symbol at a time), it encrypts plaintext in fixed-size blocks of length \(m\). Each ciphertext symbol in a block depends on all plaintext symbols of that block.

Let the plaintext block be a row vector
\[
P = [P_1\; P_2\; \cdots\; P_m]
\]
and let \(K\) be an \(m \times m\) key matrix over \(Z_{26}\).
Encryption is performed as:
\[
C = P K \pmod{26}
\]
where \(C = [C_1\; C_2\; \cdots\; C_m]\).

For decryption:
\[
P = C K^{-1} \pmod{26}.
\]
So the key matrix must be invertible modulo 26, i.e., \(\gcd(\det(K), 26)=1\).
If this condition is not satisfied, decryption is not possible.

Hill cipher resists single-letter frequency analysis due to block mixing, but is vulnerable to known-plaintext attacks.

\section{Algorithm}
The algorithm for encryption using Hill cipher is,
\begin{enumerate}
  \item Inout plaintext string \(P\).
  \item Initialize an empty ciphertext string \(C\).
  \item Process \(P\) in blocks of 2 characters:
  \begin{enumerate}
    \item Convert first character to number: \(p_1 = P[i]-'a'\).
    \item If second character exists, set \(p_2 = P[i+1]-'a'\); otherwise use padding \(x\), i.e., \(p_2 = 'x'-'a'\).
    \item Compute:
    \[
      c_1 = (k_{11}p_1 + k_{12}p_2)\bmod 26,\quad
      c_2 = (k_{21}p_1 + k_{22}p_2)\bmod 26
    \]
    \item Convert \(c_1, c_2\) back to letters and append to \(C\).
  \end{enumerate}
  \item Return ciphertext \(C\).
\end{enumerate}
The algorithm for decryption using Hill cipher is,
\begin{enumerate}
  \item Input ciphertext string \(C\).
  \item Initialize an empty plaintext string \(P\).
  \item Compute inverse key matrix \(K^{-1}\) modulo 26 (only possible if \(\gcd(\det(K),26)=1\)).
  \item Process \(C\) in blocks of 2 characters:
  \begin{enumerate}
    \item Convert block characters to numbers:
    \[
      c_1 = C[i]-'a',\quad c_2 = C[i+1]-'a'
    \]
    \item Compute:
    \[
      p_1 = (k^{-1}_{11}c_1 + k^{-1}_{12}c_2)\bmod 26,\quad
      p_2 = (k^{-1}_{21}c_1 + k^{-1}_{22}c_2)\bmod 26
    \]
    \item If any value is negative, add 26 to normalize into \([0,25]\).
    \item Convert \(p_1, p_2\) back to letters and append to \(P\).
  \end{enumerate}
  \item Return plaintext \(P\).
\end{enumerate}

% \section{Flowchart (motherboard)}
% \begin{figure}[H]
%   \centering
%   \begin{autofitflowchart}
%     \node (start) [terminator] {Start};
%     \node (input) [io, below=of start] {Input};
%     \node (process1) [process, below=of input] {Process};
%     \node (decision1) [decision, below=of process1, yshift=-0.2cm] {Condition?};
%     \node (loopstep) [process, right=of decision1, xshift=1.2cm] {Update / Loop};
%     \node (output) [io, below=of decision1, yshift=-0.3cm] {Output};
%     \node (end) [terminator, below=of output] {End};

%     \draw [line] (start) -- (input);
%     \draw [line] (input) -- (process1);
%     \draw [line] (process1) -- (decision1);

%     \draw [line] (decision1.east) -- node[above]{No} (loopstep.west);
%     \draw [line] (loopstep.north) |- ([xshift=1.0cm]process1.east) |- (process1.east);

%     \draw [line] (decision1) -- node[right]{Yes} (output);
%     \draw [line] (output) -- (end);
%   \end{autofitflowchart}
%   \caption{Flowchart title}
% \end{figure}

\section{Code}
\begin{minted}{cpp}
#include<bits/stdc++.h>
using namespace std;

class hill {
  private:
    vector<vector<int>> key;
    vector<vector<int>> inverseKey;

    public:
    hill();
    void printKey();
    string encrypt(string plaintext);
    string decrypt(string ciphertext);
    int modInverse(int a, int m);
    vector<vector<int>> matrixInverse(vector<vector<int>>& mat);
};

int hill::modInverse(int a, int m) {
  a %= m;
  if (a < 0) a += m;
  for (int i = 1; i < m; i++) {
    if ((a * i) % m == 1) return i;
  }
  return -1;
}

void hill::printKey() {
  for (const auto& row : key) {
    for (int val : row) {
      cout << val << " ";
    }
    cout << endl;
  }
}

vector<vector<int>> hill::matrixInverse(vector<vector<int>>& mat) {
  int n = mat.size();
  int det = (mat[0][0] * mat[1][1] - mat[0][1] * mat[1][0]) % 26;
  if (det < 0) det += 26;
  int detInv = hill::modInverse(det, 26);
  if (detInv == -1) {
    cerr << "Error: key matrix is not invertible modulo 26 (det=" << det << ").\n";
    exit(1);
  }
  vector<vector<int>> inv(n, vector<int>(n));
  inv[0][0] = (mat[1][1] * detInv) % 26;
  inv[0][1] = (-mat[0][1] * detInv) % 26;
  inv[1][0] = (-mat[1][0] * detInv) % 26;
  inv[1][1] = (mat[0][0] * detInv) % 26;
  for (int i = 0; i < n; ++i) {
    for (int j = 0; j < n; ++j) {
      inv[i][j] %= 26;
      if (inv[i][j] < 0) inv[i][j] += 26;
    }
  }
  return inv;
}

hill::hill(){
  srand(time(0));
  bool validKey = false;
  while (!validKey) {
    key = {{rand() % 26, rand() % 26}, {rand() % 26, rand() % 26}};
    // while(__gcd(key[0][0], 26) != 1) key[0][0] = rand() % 26;
    // while(__gcd(key[0][1], 26) != 1) key[0][1] = rand() % 26;
    // while(__gcd(key[1][0], 26) != 1) key[1][0] = rand() % 26;
    // while(__gcd(key[1][1], 26) != 1) key[1][1] = rand() % 26;
    int det = (key[0][0] * key[1][1] - key[0][1] * key[1][0]) % 26;
    if (det < 0) det += 26;
    if (modInverse(det, 26) != -1) {
      validKey = true;
    }
  }
  inverseKey = matrixInverse(key);
}

string hill::encrypt(string plaintext) {
  string ciphertext = "";
  for (int i = 0; i < plaintext.length(); i += 2) {
    int p1 = plaintext[i] - 'a';
    int p2 = (i + 1 < plaintext.length()) ? plaintext[i + 1] - 'a' : 'x' - 'a';
    int c1 = (key[0][0] * p1 + key[0][1] * p2) % 26;
    int c2 = (key[1][0] * p1 + key[1][1] * p2) % 26;
    ciphertext += (char)('a' + c1);
    ciphertext += (char)('a' + c2);
  }
  return ciphertext;
}

string hill::decrypt(string ciphertext) {
  string plaintext = "";
  for (int i = 0; i < ciphertext.length(); i += 2) {
    int c1 = ciphertext[i] - 'a';
    int c2 = ciphertext[i + 1] - 'a';
    int p1 = (inverseKey[0][0] * c1 + inverseKey[0][1] * c2) % 26;
    int p2 = (inverseKey[1][0] * c1 + inverseKey[1][1] * c2) % 26;
    if (p1 < 0) p1 += 26;
    if (p2 < 0) p2 += 26;
    plaintext += (char)('a' + p1);
    plaintext += (char)('a' + p2);
  }
  return plaintext;
}

int main() {
  hill cipher;
  cout << "Key Matrix:\n";
  cipher.printKey();
  cout << "Enter Plaintext : ";
  string text;
  getline(cin, text);
  string encrypted = cipher.encrypt(text);
  cout << "Encrypted text: " << encrypted << endl;
  string decrypted = cipher.decrypt(encrypted);
  cout << "Decrypted text: " << decrypted << endl;
  return 0;
}
\end{minted}

\newpage

\section{Result and Discussion - I/O}
While some words is encrypted and decrypted correctly, some words are not decrypted back to the original plaintext.
\begin{minted}{text}
Key Matrix:
20 19
25 2
Enter Plaintext : password
Encrypted text: olaseghp
Decrypted text: password
\end{minted}
Others have some random letters in the decrypted text, which is due to the padding character 'x' used when the plaintext length is odd.
\begin{minted}{text}
Key Matrix:
25 6
0 5
Enter Plaintext : encrypt
Encrypted text: wnwhoxpl
Decrypted text: encryptx
\end{minted}


\section{Conclusion}
Hill cipher demonstrates how linear algebra applies to cryptography, but its small block size limits practical security. Larger matrices improve security at the cost of complexity.
